% Part: computability
% Chapter: recursive-functions
% Section: notations-pr-functions

\documentclass[../../../include/open-logic-section]{subfiles}

\begin{document}

\olfileid{cmp}{rec}{not}
\olsection{આદિમ પુનરાવર્તી વિધેયોનાં સંકેતન}

આદિમ પુનરાવર્તી વિધેયોનું ચોક્કસ આગમનાત્મક વર્ણન આપ્યું છે તેનો એક
ફાયદો એ છે કે આપણે તેમને વ્યવસ્થિત રીતે દર્શાવી શકીએ છીએ.
દાખલા તરીકે, આવા દરેક વિધેયને નીચે મુજબ એક ``સંકેતન'' આપી શકીએ.
શૂન્ય વિધેય, ઉત્તરગામી વિધેય અને પ્રક્ષેપ વિધેયો માટે અનુક્રમે
$\Zero$, $\Succ$ અને $\Proj{n}{i}$ સંકેતો વાપરીએ. હવે ધારો કે
$k$-આર્ગ્યુમેન્ટવાળું વિધેય~$f$ અને $n$-આર્ગ્યુમેન્ટવાળાં વિધેયો
$g_0$, \dots,~$g_{k-1}$ના સંયોજનથી $h$ વ્યાખ્યાયિત છે, અને આ
વિધેયોને અનુક્રમે $F$, $G_0$, \dots,~$G_{k-1}$ સંકેતનો આપ્યાં છે.
તો નવો સંકેત $\fn{Comp}_{k,n}$ વાપરીને $h$ વિધેયને
$\fn{Comp}_{k,n}[F,G_0,\dots,G_{k-1}]$ વડે દર્શાવી શકીએ.

આદિમ પુનરાવર્તનથી વ્યાખ્યાયિત વિધેયો માટે પણ આવાં જ સંકેતનો
વાપરી શકીએ. ધારો કે $(k+1)$-આર્ગ્યુમેન્ટવાળું વિધેય~$h$,
$k$-આર્ગ્યુમેન્ટવાળા વિધેય~$f$ અને $(k+2)$-આર્ગ્યુમેન્ટવાળા
વિધેય~$g$માંથી આદિમ પુનરાવર્તન દ્વારા વ્યાખ્યાયિત છે; અને $f$ તથા
$g$ને અપાયેલાં સંકેતનો અનુક્રમે $F$ અને~$G$ છે. ત્યારે~$h$ને
અપાતું સંકેતન $\fn{Rec}_k[F,G]$ છે.

યાદ કરો કે સરવાળાનું વિધેય આદિમ પુનરાવર્તનથી આમ વ્યાખ્યાયિત છે:
\begin{align*}
  \Add(x_0, 0) & = \Proj{1}{0}(x_0) = x_0\\
  \Add(x_0, y+1) & = \Succ(\Proj{3}{2}(x_0, y, \Add(x_0, y))) = \Add(x_0, y) +1
\end{align*}
અહીં $f$નું સ્થાન $\Proj{1}{0}$ લે છે અને $g$નું સ્થાન
$\Succ(\Proj{3}{2}(x_0, y, z))$ લે છે. પછીના વિધેયને
$\fn{Comp}_{1,3}[\Succ,\Proj{3}{2}]$ સંકેતન મળે છે, કારણ કે તે
$1$-આર્ગ્યુમેન્ટવાળા વિધેય~$\Succ$ અને
$3$-આર્ગ્યુમેન્ટવાળા વિધેય~$\Proj{3}{2}$ના સંયોજનથી વ્યાખ્યાયિત
થાય છે. આ રીતે સરવાળાનું વિધેય નીચેના સંકેતનથી દર્શાવી શકીએ:
\[
\fn{Rec}_1[\Proj{1}{0},\fn{Comp}_{1,3}[\Succ,\Proj{3}{2}]].
\]
આવાં સંકેતનો ક્યારેક ઉપયોગી થાય છે; દાખલા તરીકે, આદિમ પુનરાવર્તી
વિધેયોની પરિગણના કરતી વખતે.

\begin{prob}
  $\Mult$ માટે સંપૂર્ણ આદિમ પુનરાવર્તી સંકેતન આપો.
\end{prob}

\end{document}
